\DocumentMetadata{
  pdfversion=2.0,
  pdfstandard=ua-2,
  lang=en-US,
  tagging=on,
  tagging-setup={math/setup=mathml-SE,tabsorder=structure}
}
\documentclass[11pt,letterpaper]{article}
\usepackage[margin=0.68in,includefoot]{geometry}
\usepackage{newcomputermodern}
\usepackage{amsmath,amscd,mathtools}
\providecommand{\square}{\mdlgwhtsquare}
\usepackage[hidelinks]{hyperref}
\hypersetup{
  pdftitle={Combinatorics PhD Exam, August 2015},
  pdfauthor={University of Florida Department of Mathematics},
  pdfsubject={Graduate mathematics examination},
  pdfdisplaydoctitle=true,
  bookmarksopen=true
}
\setcounter{secnumdepth}{0}
\setlength{\parindent}{0pt}
\setlength{\parskip}{0.28em}
\setlength{\emergencystretch}{3em}
\setlength{\leftmargini}{2.15em}
\setlength{\leftmarginii}{2.65em}
\setlength{\leftmarginiii}{3em}
\renewcommand{\labelenumi}{\textbf{\theenumi.}}
\pagestyle{empty}
\raggedbottom
\allowdisplaybreaks
\newenvironment{examproblems}[1]{%
  \begin{enumerate}%
  \setcounter{enumi}{#1}%
  \setlength{\topsep}{0.35em}%
  \setlength{\partopsep}{0pt}%
  \setlength{\itemsep}{0.48em}%
  \setlength{\parsep}{0.18em}%
}{\end{enumerate}}
\newenvironment{examsubparts}{%
  \begin{enumerate}%
  \setlength{\topsep}{0.18em}%
  \setlength{\partopsep}{0pt}%
  \setlength{\itemsep}{0.16em}%
  \setlength{\parsep}{0.1em}%
}{\end{enumerate}}
\newenvironment{examinstructions}{%
  \par\begingroup\itshape
}{%
  \par\endgroup\vspace{0.18em}
}
\makeatletter
\renewcommand\section{\@startsection{section}{1}{0pt}%
  {-1.25ex plus -.25ex minus -.1ex}%
  {0.55ex plus .1ex}%
  {\normalfont\large\bfseries}}
\renewcommand{\@maketitle}{%
  \newpage\null\vskip -1.2em%
  {\footnotesize\bfseries
    \href{https://www.ufl.edu/}{University of Florida}, \href{https://math.ufl.edu/}{Department of Mathematics}\par}%
  \vskip 0.18em%
  \tagstructbegin{tag=Title}%
    {\LARGE\bfseries\href{https://gma.math.ufl.edu/exams/combinatorics-phd/}{Combinatorics PhD Exam}, August 2015\par}%
  \tagstructend%
  \vskip 0.35em%
  \tagmcbegin{artifact}%
    \hrule height 1pt%
  \tagmcend%
  \vskip 0.55em%
}
\makeatother
\title{Combinatorics PhD Exam, August 2015}
\author{}
\date{}
\begin{document}
\maketitle
\thispagestyle{empty}
\begin{examproblems}{0}
\item
Let \(f(n)\) be the number of permutations of length~\(n\) that contain no cycles of length exactly two.
\begin{examsubparts}
\item[\textnormal{(a)}]
Find the exponential generating function of the sequence \(f(n)\). You can assume that \(f(0)=1\).
\item[\textnormal{(b)}]
Find \(\lim_{n\to\infty} f(n)/n!\).
\end{examsubparts}
\item
We select an odd number of people from a group of~\(n\) people to form a committee. Then we select an even number of people from this committee to serve on a subcommittee. (Zero is an even number.) In how many ways can we do this?
\item
Let~\(T\) be a tree in which no path is longer than~\(k\). Is it true that there is a vertex \(v\in T\) that is contained in all paths of length~\(k\) in~\(T\)?
\item
Let \(X_n(S)\) be the size of the subset~\(S\) of \([n]\) selected uniformly, that is, each element of \([n]\) has fifty percent chance to be selected. Compute \(\mathrm{VAR}(X_n)\).
\item
Let \(f_n\) be the number of all parts of all compositions of~\(n\). So \(f_1=1\), \(f_2=3\), and \(f_3=8\). Find and prove an explicit formula for \(f_n\).
\item
Let \(m_n\) be the number of ways to have a group of~\(n\) people split into nonempty subsets, to have each subset sit down around a circular table, then to arrange the tables in a circle. Find the exponential growth rate of the sequence \(m_n/n!\). Two arrangements are considered identical if each person has the same left neighbor in them, and each table has the same left neighbor in them.
\item
Prove that every convex polyhedron has two faces that have the same number of vertices.
\item
Let~\(p\) be a permutation that can be decomposed into the disjoint union of~\(k\) increasing subsequences, but not into the disjoint union of \(k-1\) or fewer increasing subsequences. What can be said about the length of the longest decreasing subsequence of~\(p\)?
\end{examproblems}
\end{document}
